% ---------------------------------------------------------------------
% Drop-in: the residual bound, in norm and then at every instant.
%
% Structure follows the shortest route.  The norm statement is one line
% -- the fit cannot be farther from the data than the nominal is, and
% the nominal leaves e + n -- and that is (VIII.3).  The pointwise
% statement is the same line plus exactly one term, which is what the
% naive version omits and why it fails at the onset.
%
% No triangle inequality on the fit displacement appears anywhere:
% r = (I-P)y is an identity, and splitting it gives the pointwise form
% directly.  Only two inequalities are used, both exact in origin:
% ||(I-P)v|| <= ||v||, and |v_i| <= sqrt(P_ii) ||v|| for v in the span.
%
% An earlier version routed the same result through the projector's
% kernel and a box argument on the band.  This one assumes less (the
% linear case needs no hypothesis at all, since (P.6) is an identity)
% and measures tighter, worst ratio 0.611 against 0.763.
%
% Constants regenerated by
%     python analysis/failing_runs.py        DataSet/exp
%     python analysis/projection_constants.py
%
% Equation tags are local (P.1)-(P.8); renumber to the host document.
% ---------------------------------------------------------------------

\documentclass[10pt]{article}
\usepackage[a4paper,margin=22mm]{geometry}
\usepackage{amsmath,amssymb,booktabs}
\setlength{\parskip}{0.35em}\setlength{\parindent}{0pt}
\newcommand{\Mdot}{\dot M}
\newcommand{\rbar}{\bar\rho}
\title{\vspace{-16mm}\textbf{The residual, in norm and at every
instant}\vspace{-3mm}}
\date{\small drop-in for Sec.~VIII; the norm form is (VIII.3)}
\begin{document}\maketitle

\subsection{The statement in norm, in one line}

Write the recorded rate over the post-onset window as the nominal
response, the deviation the model does not carry, and sensor noise,
\begin{equation}
  y = \omega_{\rm nom} + e_\omega + n ,
  \tag{P.1}
\end{equation}
and let $\hat\omega$ minimise $\|g-y\|$ over the model family $M$. The
nominal is itself a member of $M$ --- it is a hyperbolic cosine with the
true amplitude and exponent --- so it is one of the candidates the
minimiser was choosing among, and therefore
\begin{equation}
  \|r\| \;=\; \bigl\|\hat\omega - y\bigr\|
  \;=\; \min_{g\in M}\|g-y\|
  \;\le\; \bigl\|\omega_{\rm nom} - y\bigr\|
  \;=\; \|e_\omega + n\| .
  \tag{P.2}
\end{equation}
That is the whole argument in norm. It uses optimality and nothing else:
no linearity, no projection, no expansion, no triangle inequality. It is
(VIII.3), and dividing by $\sqrt{\tau_{\rm end}}$ puts it in RMS form,
$\mathrm{RMS}(r) \le \mathrm{RMS}(e_\omega) + \mathrm{RMS}(n)$.

\subsection{The same statement at every instant}

The pointwise version one would like to write is
\begin{equation}
  |r_i| \;\overset{?}{\le}\; |e_{\omega,i}| + k\sigma_n ,
  \tag{P.3}
\end{equation}
and it does not follow from (P.2). A norm inequality constrains a sum of
$N$ numbers; about any one of them it says only $|r_i|\le\|r\|$, which
is a bound by the total and not by the local value $|e_{\omega,i}|$.

Nor is (P.3) true, and where it fails is the onset. With $C_2$ held at
the value the rig fixes, the admissible curves form a two-dimensional
subspace
\begin{equation}
  V = \operatorname{span}\{u,\mathbf 1\},
  \qquad u_i = \cosh C_2\tau_i - 1 ,
  \tag{P.4}
\end{equation}
the least-squares fit is $\hat\omega = Py$ with $P$ the orthogonal
projector onto $V$, and $P$ has the closed form
\begin{equation}
  P_{ij} \;=\; \frac1N \;+\;
    \frac{\tilde u_i\tilde u_j}{\sum_k\tilde u_k^{\,2}},
  \qquad \tilde u = u - \bar u .
  \tag{P.5}
\end{equation}
The leading $1/N$ alone settles it: a perturbation anywhere in the
window shifts the fitted \emph{constant}, and so moves the fitted curve
at every sample including the first. Since $e_\omega$ lives at the far
end of the window --- its envelope is $E\propto\sinh C_2\tau$ --- its
mean displaces the fit at $\tau = 0$ by roughly $0.3\sup E$ while
$e_\omega(0)$ itself vanishes. No constant $\Gamma$ rescues a
multiplicative form either: measured on the campaign the ratio of the
true bound to $E$ is $2.1$ at the end of the window and $27$ to $95$ at
the first sample.

What (P.3) omits is exactly that displacement, and writing it down costs
one term. Since $\hat\omega = Py$ is an identity --- true whether or not
the nominal lies in $V$ --- substituting (P.1) gives
\begin{equation}
  r \;=\; (I-P)\,y
    \;=\; \underbrace{(I-P)\,\omega_{\rm nom}}_{\text{misspecification}}
      \;+\; \bigl(e_\omega + n\bigr)
      \;-\; \underbrace{P\bigl(e_\omega + n\bigr)}_{\text{displacement}} ,
  \tag{P.6}
\end{equation}
exactly, and with no hypothesis at all. Two remarks close it.

\emph{The middle group is (P.3)'s right-hand side.} Taking absolute
values sample by sample returns $|e_{\omega,i}| + |n_i|$, which is what
(P.3) already had.

\emph{The last group is bounded because it lives in two dimensions.}
$P(e_\omega+n)\in V$, and for $v\in V$ we have $v = Pv$, so
$v_i = \langle P_{i\cdot},v\rangle$ with
$\|P_{i\cdot}\|^2 = (P^2)_{ii} = P_{ii}$, whence
$|v_i| \le \sqrt{P_{ii}}\,\|v\|$ --- an identity followed by
Cauchy--Schwarz, not an estimate. And $\|P(e_\omega+n)\| \le
\|e_\omega+n\|$ because $P$ is an orthogonal projection, so (P.2)
bounds it too.

Together,
\begin{equation}
  \boxed{\;
  |r_i| \;\le\; \bigl|\bigl[(I-P)\omega_{\rm nom}\bigr]_i\bigr|
  \;+\; |e_{\omega,i}| \;+\; |n_i|
  \;+\; \sqrt{P_{ii}}\,\bigl\|e_\omega + n\bigr\| \;}
  \tag{P.7}
\end{equation}
which is (P.3) with two additions: the displacement, the price of the
fit being free to move, and the misspecification, the price of the
calibrated exponent not being the true one. The first is unavoidable;
the second vanishes when the exponents agree.

\paragraph{The noise, separated.} Splitting
$\|e_\omega+n\| \le \|e_\omega\| + \|n\|$ and using
$|e_{\omega,i}|\le E_i$, $\|e_\omega\|\le\|E\|$ from (93) leaves the
noise in two places, $|n_i|$ and $\sqrt{P_{ii}}\|n\|$. The second must
not be written $\sigma_n\sqrt N$: for $n\sim\mathcal N(0,\sigma_n^2I)$
and $P$ symmetric idempotent,
$\mathbb E\|Pn\|^2 = \sigma_n^2\operatorname{tr}P = 2\sigma_n^2$
whatever $N$ is, and $\operatorname{Var}[(Pn)_i] = \sigma_n^2P_{ii}$, so
the two together are at most $k\sigma_n(1+\sqrt{P_{ii}})$ --- against a
factor $\sqrt{N/2} = 4.5$ at $N = 41$ for the crude route. The form
checked below is therefore
\begin{equation}
  |r_i| \;\le\; \underbrace{m_i}_{\text{misspec.}}
  \;+\; \underbrace{E_i}_{\text{pointwise}}
  \;+\; \underbrace{\sqrt{P_{ii}}\,\|E\|}_{\text{displacement}}
  \;+\; k\sigma_n\bigl(1+\sqrt{P_{ii}}\bigr),
  \quad m = \bigl|(I-P)\omega_{\rm nom}\bigr| .
  \tag{P.8}
\end{equation}
Two Gaussian events enter per sample, so a union bound over the window
replaces $k$ by $\Phi^{-1}(1-\alpha/4N)$: for $N = 53$ that is $3.50$ at
$95\%$ and $3.90$ at $99\%$, against $1.96$ and $2.58$ for one sample.

\paragraph{If $C_2$ is fitted rather than given.} There is then no $P$
and (P.6) is unavailable --- but (P.2) is untouched, having never used
one. Only the pointwise step needs replacing, by the corresponding
inequality on the two-dimensional manifold, where the displacement
carries a factor $\bigl(1-\|e_\omega+n\|/\varrho\bigr)^{-1}$ with
$\varrho$ its reach. That is the calibration stage's case. The per-run
fit, which is handed a $C_2$, is the linear one throughout.

\subsection{The constants, and why they do not grow}

Everything in (P.8) is computable before a run is flown: $P$ depends
only on $C_2$ and the sample grid, $E$ only on the geometry box.

\begin{center}\small
\begin{tabular}{@{}rrrrrr@{}}
\toprule
$\Mdot$ & $N$ & $\operatorname{tr}P$ & $\kappa=\max_i\sqrt{P_{ii}}$ &
$\|E\|/\sup E$ & $\kappa\|E\|/\sup E$\\
N\,m/s & & & & & \\
\midrule
$0.10$ & $74$ & $2.000$ & $0.352$ & $3.13$ & $1.12$\\
$0.20$ & $63$ & $2.000$ & $0.362$ & $3.12$ & $1.15$\\
$0.30$ & $57$ & $2.000$ & $0.370$ & $3.10$ & $1.17$\\
$0.45$ & $54$ & $2.000$ & $0.376$ & $3.09$ & $1.18$\\
$0.65$ & $48$ & $2.000$ & $0.385$ & $3.07$ & $1.21$\\
$0.90$ & $44$ & $2.000$ & $0.395$ & $3.02$ & $1.22$\\
$1.20$ & $41$ & $2.000$ & $0.404$ & $2.99$ & $1.23$\\
\bottomrule
\end{tabular}
\end{center}

The trace is $2$ at every rate and every window length, and (P.5) shows
why directly: $\sum_i 1/N = 1$ and
$\sum_i\tilde u_i^2/\sum_k\tilde u_k^2 = 1$, one unit for each basis
direction. That is why the constants do not grow with the record --- a
two-dimensional projector cannot move much, however many samples it is
given. The displacement term is $1.09$ to $1.29$ times $\sup E$ across
all $140$ runs, so in uniform form
$\max_i|r_i| \lesssim 2.3\sup E + 1.4\,k\sigma_n$.

Equation (P.5) also locates the worst sample: $P_{ii}$ is largest where
$|\tilde u_i|$ is, which is the last sample of the window in all $140$
runs, and only $1/N = 0.013$ to $0.024$ of $P_{ii}$ there is the
constant --- the rest is the amplitude direction, whose $\cosh$ has
grown by then. It reproduces the numerically formed projector to
$10^{-16}$.

\subsection{The bound tested, and the two things it exposes}

Applied to every sample of every run at $k = 3.5$:

\begin{center}\small
\begin{tabular}{@{}rrrrrr@{}}
\toprule
$\Mdot$ & \multicolumn{2}{c}{runs entirely inside}
        & worst & deployed & $\sigma_{\rm pre}$ / $\sigma_{\rm in}$\\
N\,m/s & $\sigma_{\rm pre}$ & $\sigma_{\rm in}$ & ratio & worst
       & $^\circ$/s\\
\midrule
$0.10$ & $20/20$ & $20/20$ & $0.324$ & $0.369$ & $0.087$ / $0.831$\\
$0.20$ & $20/20$ & $20/20$ & $0.391$ & $0.551$ & $0.044$ / $0.800$\\
$0.30$ & $18/20$ & $20/20$ & $0.457$ & $0.584$ & $0.101$ / $0.823$\\
$0.45$ & $19/20$ & $20/20$ & $0.491$ & $0.979$ & $0.127$ / $0.707$\\
$0.65$ & $18/20$ & $20/20$ & $0.580$ & $0.863$ & $0.122$ / $0.856$\\
$0.90$ & $15/20$ & $20/20$ & $0.556$ & $1.030$ & $0.111$ / $0.749$\\
$1.20$ & $13/20$ & $20/20$ & $0.611$ & $1.212$ & $0.173$ / $0.773$\\
\bottomrule
\end{tabular}
\end{center}

The bound holds on all $140$ runs at every sample, the worst reaching
$0.611$ of it. For comparison, dropping the displacement term --- that
is, checking (P.3) itself with $E$ in place of $|e_\omega|$ --- holds
for $99.5\%$ of samples and $120$ of $140$ runs entirely. The naive form
is nearly right, and the extra term is what closes the last twenty runs.

\paragraph{The noise scale must be measured in the window.} Taking
$\sigma_n$ from the pre-onset record leaves $123$ of $140$ inside, and
the shortfall is at the fast end. That is not the bound failing but its
hypothesis: before the onset the vehicle is carried by every landing
gear, after it pivots on one edge, and the in-window disturbance is
$0.71$ to $0.86^\circ$/s against a pre-onset floor of $0.04$ to $0.17$.
The honest estimate is the in-window content above $5$~Hz, which the
fitted model provably cannot produce --- $\hat\omega$ is smooth on the
scale of the window --- so whatever is there is disturbance by
construction rather than by assumption.

\paragraph{A pinned amplitude adds a term.} $V$ is the family the bound
is about: amplitude and level free. The shipped estimator frees neither.
It takes $C_1 = K\Mdot$ from the calibration stage, and it forms the
level $C$ from the \emph{pre-onset} segment rather than minimising over
it on the window, so once the onset is chosen its post-onset curve has
no free parameter and $\hat\omega\ne Py$. Two terms then survive,
\begin{equation*}
  |r_i| \;\le\; |\Delta C_1|\,u_i + |C-b|
        + m_i + E_i + \sqrt{P_{ii}}\,\|E\|
        + k\sigma_n\bigl(1+\sqrt{P_{ii}}\bigr),
\end{equation*}
the first being the calibration's amplitude error and the second the
baseline channel of (105)--(106). Their measured sizes are very
different. The amplitude term is $|\Delta C_1|/C_1 = 11.3\%$ and it is
what puts $26$ runs outside the norm form (VIII.3); without it the
deployed residual reaches $1.21$ of the bound. The baseline term is
$0.005$ to $0.029^\circ$/s, under $2\%$ of $\sup E$. One constraint
matters and the other does not, and the bound says which.

\subsection{What the misspecification term is worth}
\label{sec:c2}

The term $m = |(I-P)\omega_{\rm nom}|$ is the whole content of the
objection that the fitted curve and the true one do not share a basis.
They do not: the nominal carries $C_2^{*}$ and the span is built on the
calibrated $C_2$. Writing $C_2 = C_2^{*}(1+\varepsilon_2)$, it is the
price of conditioning on the calibration stage --- which fits $C_2$
freely, by a per-run nonlinear least squares and a per-configuration
median --- and it is small, almost linear in $\varepsilon_2$:

\begin{center}\small
\begin{tabular}{@{}rrrrr@{}}
\toprule
$\Mdot$ & \multicolumn{4}{c}{$\bigl\|(I-P)\omega_{\rm nom}\bigr\|_\infty
 \,/\,\sup E$}\\
N\,m/s & $\varepsilon_2 = 1\%$ & $2\%$ & $5\%$ & $10\%$\\
\midrule
$0.10$ & $0.005$ & $0.010$ & $0.027$ & $0.059$\\
$0.45$ & $0.013$ & $0.027$ & $0.071$ & $0.152$\\
$0.90$ & $0.018$ & $0.036$ & $0.093$ & $0.195$\\
$1.20$ & $0.020$ & $0.041$ & $0.105$ & $0.221$\\
\bottomrule
\end{tabular}
\end{center}

A $2\%$ exponent error costs $1$ to $4\%$ of the envelope and a $5\%$
error costs $3$ to $11\%$, the fast ramp being the more exposed because
its window is shorter and the two hyperbolic shapes are less separated
there. Read the other way, $98\%$ of a nominal misspecified by $1\%$
still lies inside $V$ --- the same near-collinearity of $\cosh$ and
$\sinh$ over this window, $3.4^\circ$ once the constant is removed, that
Sec.~VI-E uses elsewhere. A shifted exponent barely leaves the span it
is fitted in.

Like the pinned amplitude the term enters additively and is not covered
by $\rbar$, which bounds the physics the model omits and says nothing
about the constants the model is given. The empirical check above is
unaffected by any of it: it compares a \emph{measured} record, which
carries whatever exponent the rig actually has, against the envelope, so
the $140/140$ already includes the misspecification. What the term
supplies is the proof that this had to be so, and a budget line for it.

\end{document}
